% 02_c07_heat_treatment_scales
\documentclass[tikz,border=8pt]{standalone}
\usepackage{xfab-tikz}
\begin{document}
\begin{tikzpicture}[x=1cm,y=1cm]
% temperature axis 0..1300 C mapped to 0..10.4 cm
\def\T#1{#1*0.008}
\draw[cToolD, line width=0.8pt] (0,0) -- (10.6,0);
\foreach \t in {0,200,...,1200}{\draw[cToolD] (\T{\t},0) -- (\T{\t},-0.1); \node[font=\footnotesize, anchor=north] at (\T{\t},-0.1) {\t};}
\node[lab, anchor=north] at (5.2,-0.55) {temperature (°C)};
\foreach \y/\a/\b/\name/\what/\time/\c in {
  6.0/60/100/boiled egg/proteins unfold and bond/minutes/cFeed,
  5.1/120/180/{curing of thermosets, composites}/chains cross-link/hours/cChem,
  4.2/150/650/tempering of steel/{carbides precipitate, stress relaxes}/hours/cIPart,
  3.3/800/900/quench hardening of steel/{austenite, then martensite on quenching}/seconds/cHeat,
  2.4/600/900/stress relief of AM metal parts/dislocations recover/hours/cIPart,
  1.5/900/1250/hot isostatic pressing (HIP)/{creep and diffusion close pores, 100--200 MPa}/hours/cHeat}{
  \fill[fill=\c!60, draw=\c!60!black, line width=0.5pt, rounded corners=2pt] (\T{\a},\y-0.22) rectangle (\T{\b},\y+0.22);
  \node[lab, font=\small, anchor=south west] at (0,\y+0.25) {\textbf{\name}\enspace{\color{inksoft}\what\ · \time}};}
\draw[steellight, dashed] (\T{20},0) -- (\T{20},6.6);
\node[lab, text=inksoft, anchor=south] at (\T{20},6.6) {room};
\xfabref{0}{-1.15}{Typical ranges, after Ashby, Shercliff \& Cebon, \emph{Materials}, 3rd ed., 2013, ch.~19, and ASTM F3301 (heat treatment of AM metal parts).}

\end{tikzpicture}
\end{document}
